\documentclass[11pt,UTF8,fontset=fandol]{ctexart}
\usepackage{amsmath,amssymb,amsfonts}
\usepackage{graphicx,float,caption,subcaption}
\usepackage{booktabs,array,multirow,makecell}
\usepackage{geometry,enumitem,xcolor}
\usepackage[hidelinks]{hyperref}
\usepackage{url}
\geometry{margin=2.3cm}
\graphicspath{{/Users/yye/project/论文/assets/higpt/}}
\captionsetup{font=small,labelfont=bf}
\setlength{\parskip}{0.4em}
\hypersetup{pdftitle={HiGPT 中文精读导读},pdfauthor={Tang et al.}}

% 小白知识卡片：上下加线 + 粗体标签，稳健不挑内容
\newenvironment{tip}[1]{%
  \par\smallskip\noindent\rule{\linewidth}{0.4pt}\par\nobreak
  \noindent\textbf{【小白补丁】#1}\par\nobreak\noindent\ignorespaces
}{%
  \par\noindent\rule{\linewidth}{0.4pt}\par\smallskip
}

\title{HiGPT 中文精读导读：教大语言模型读懂“异质图”}
\author{原文：Tang, Yang, Wei, Shi, Xia, Yin, Huang（KDD 2024，香港大学 \& 百度）}
\date{}

\begin{document}
\maketitle
\begin{center}\small\textit{本导读为面向初学者的中文学习笔记：含要点翻译、公式精讲与图表解读。}\end{center}

\section*{论文卡片}
\begin{center}
\begin{tabular}{@{}ll@{}}
\toprule
\textbf{标题} & HiGPT: Heterogeneous Graph Language Model \\
            & （HiGPT：异质图语言模型） \\
\textbf{作者} & Jiabin Tang, Yuhao Yang, Wei Wei, Lei Shi, Long Xia, Dawei Yin, Chao Huang \\
\textbf{机构} & 香港大学（University of Hong Kong）、百度（Baidu Inc.） \\
\textbf{会议} & KDD 2024（数据挖掘领域顶级会议） \\
\textbf{arXiv} & 2402.16024 \\
\textbf{代码} & \url{https://github.com/HKUDS/HiGPT} \\
\textbf{主页} & \url{https://HiGPT-HKU.github.io} \\
\textbf{方向} & 异质图（Heterogeneous Graph）$\times$ 大模型 $\times$ 图指令微调 \\
\textbf{一句话} & 把“图指令微调”推广到多类型节点/边的\emph{异质图}，做一个不用在下游 \\
              & 数据上微调、就能跨数据集迁移的“通用异质图大模型”。 \\
\bottomrule
\end{tabular}
\end{center}

\section{一、这篇论文在解决什么问题}

现实中的很多数据天生是“图”：推荐系统里有“用户—商品”，知识图谱里有“实体—关系”，
学术网络里有“论文—作者—会议”。这些图有一个共同特点：\textbf{节点和边不止一种类型}。
比如学术网络里，“论文”“作者”“会议”是三种不同的节点，“论文被作者写”“论文发在会议上”
是不同的边。这种\textbf{多类型节点、多类型边}的图，就叫\textbf{异质图（heterogeneous graph）}；
反过来，所有节点都同一类、所有边都同一类的图叫\textbf{同质图（homogeneous graph）}。

为异质图设计的神经网络叫 HGNN（异质图神经网络，Heterogeneous Graph Neural Network），
它们已经做得很好。但作者指出一个根本痛点：\textbf{现有 HGNN 几乎都是“预训练 + 微调”在\emph{同一个}
数据集上跑}——在数据集 A 上训练好的模型，几乎\emph{没法}直接搬到数据集 B 上用。原因有三：

\begin{enumerate}[leftmargin=2em]
  \item \textbf{关系类型异质性漂移（Relation Type Heterogeneity Shift）}：不同图的“关系种类”完全不同。
        推荐图里的边是“点击/收藏/评论/购买”，学术图里的边却是“论文—论文/作者—论文/论文—会议”。
        在一种关系上学好的模型，换一批全新的关系类型就懵了。这是本文要攻克的\textbf{头号难题}。
  \item \textbf{异质图结构复杂}：模型既要分辨“不同类型节点之间”的关系（异质关系），又要保留“同类型节点内部”
        的区别（同质关系），还要把这些结构信息塞进大语言模型（LLM）里去。
  \item \textbf{微调数据稀缺}：真实场景（如推荐系统的冷启动用户/商品）往往\textbf{标注极少}，
        没有足够数据去微调一个专用模型。
\end{enumerate}

\textbf{本文的目标}：做一个\textbf{通用、可迁移}的异质图大模型 HiGPT——
能从\emph{任意}异质图里学习，\textbf{不需要在下游数据集上做任何微调}，
即使遇到“节点类型集合变了、关系类型也变了”的新图也照样能用。

\begin{tip}{“关系类型异质性漂移”到底难在哪？}
打个比方：你训练了一个只认识“电影—导演—演员”关系的模型。现在换到学术图，这里只有
“论文—作者—会议”。新图里\emph{每一种关系的名字、含义、数量}都和旧图对不上——
就像一个只学过中文菜谱的厨师，突然拿到一本法文菜谱，连食材名都不认识。
传统 HGNN 把“关系种类”写死在了网络参数里，所以换图就失效。HiGPT 的核心创新，
就是让模型\textbf{在用的时候临时“现学”当前这张图有哪些关系类型}，从而适配任意新图。
\end{tip}

\section{二、摘要中译}

\begin{quote}
异质图学习的目标，是捕捉异质图中实体之间复杂的关系与多样的关系语义，从而为节点和边获得有意义的表示。
近期，异质图神经网络（HGNN）通过考虑关系异质性、使用专门的消息函数与聚合规则，取得了最先进的性能。
然而，现有的异质图学习框架在\textbf{跨不同异质图数据集泛化}方面存在局限：它们大多遵循在\emph{同一数据集}
上“预训练—微调”的范式，这限制了其适应新的、未见数据的能力。这就引出一个问题：“我们能否让异质图模型
很好地适配各种下游学习任务，即使面临节点 token 集合与关系类型异质性上的分布漂移？”为应对这些挑战，
我们提出 HiGPT，一个带\textbf{异质图指令微调（Heterogeneous graph instruction-tuning）}范式的通用大图模型。
我们的框架无需对下游数据集做任何微调，就能从任意异质图中学习。为处理异质性上的分布漂移，我们引入一个
\textbf{上下文异质图 tokenizer（in-context heterogeneous graph tokenizer）}，它能捕捉不同异质图中的语义关系，
便于模型适配。我们把大规模的、具异质性感知的图指令注入 HiGPT，使模型能有效理解复杂的关系异质性、
并区分各种类型的图 token。此外，我们引入\textbf{混合思维（Mixture-of-Thought, MoT）指令增强}范式，
通过生成多样而信息丰富的指令来缓解数据稀缺。通过在多种设定下的全面评估，我们的框架在泛化性能上表现卓越，
超越了当前领先的基准。
\end{quote}

\section{三、四大核心贡献}
\begin{enumerate}[leftmargin=2em]
  \item \textbf{提出异质图指令微调范式}：首次把“图指令微调”推广到异质图，做出一个不依赖下游微调、
        可跨数据集迁移的通用异质图大模型 HiGPT。
  \item \textbf{上下文异质图 tokenizer}：用语言来编码节点/边的类型，让模型在推理时\emph{现学}当前图的
        异质结构，从而克服“关系类型异质性漂移”。
  \item \textbf{混合思维（MoT）指令增强}：用多种提示工程（CoT/ToT/PanelGPT/GKP）让强 LLM 自动生成多样的、
        带推理过程的指令，缓解少样本场景的数据稀缺。
  \item \textbf{充分的实证验证}：在少样本、零样本、跨数据集/跨领域迁移以及“图上下文学习”等多种设定下，
        HiGPT 大幅超越各类基线（同质 GNN、异质 GNN、自监督方法）。
\end{enumerate}

\section{四、预备知识（小白先看这里）}

\begin{tip}{什么是“图指令微调”（graph instruction-tuning）}
“指令微调”原本是 LLM 领域的概念：把任务写成“一句人话的指令 + 期望回答”，喂给模型去学，
让它学会“听懂指令、照着做”。\textbf{图指令微调}就是把图也塞进这条指令里——
指令文本里夹着一个特殊的 \texttt{<graph>} 占位符，背后对应着“图编码器输出的图向量”。
模型于是学会“一边看图、一边听指令、一边回答”。本文把它从同质图推广到了异质图。
\end{tip}

\begin{tip}{什么是“元路径”与“元关系”（metapath / meta relation）}
在异质图里，描述“一条边连接了什么类型”需要一个三元组：
$\langle\,\tau(u),\ \rho(e),\ \tau(v)\,\rangle$，意思是“源节点类型—边类型—目标节点类型”，
比如 $\langle\text{电影},\ \text{导演},\ \text{导演}\rangle$，这叫\textbf{元关系（meta relation）}。
把若干元关系串起来形成的“路径模板”叫\textbf{元路径（metapath）}，
比如“作者—论文—作者（A-P-A）”表示“两位作者合写过同一篇论文”。
元路径是异质图里表达“高阶关系”的常用工具。
\end{tip}

\begin{tip}{什么是 tokenizer（在本文语境下）}
在 NLP 里，tokenizer 是“把文字切成一个个 token（词元）”的工具。本文借用这个词，
指的是\textbf{把一张异质图编码成一串“图 token”（图向量序列）}的模块——
作用类似一个图编码器（HGNN）。它的特别之处是“in-context（上下文式）”：
不把关系类型写死，而是根据输入图\emph{当场}决定怎么编码。
\end{tip}

形式化记号（看不懂可跳过，用到时再回看）。一张异质图记作
$G(V,E,A,\mathcal{T},\mathcal{R},X)$：$V$ 是节点集、$E$ 是边集、$A$ 是邻接矩阵；
$\mathcal{T}$ 是\textbf{节点类型集}、$\mathcal{R}$ 是\textbf{边（关系）类型集}（这两个集合正是“异质”的来源）；
$X=\{X_{T_i}\in\mathbb{R}^{|V_{T_i}|\times d_{T_i}}\}$ 是各类型节点的特征矩阵。
异质图神经网络（HGNN）通过“按类型分别设计的消息传递 + 聚合”来更新节点表示：
\begin{equation*}
  h_v^{(l)}=\operatorname*{Aggregate}_{\forall u\in\mathcal{N}(v),\ \forall e\in E(u,v)}
  \Big(\mathrm{Propagate}\big(h_u^{(l-1)};\,h_v^{(l-1)},e\big)\Big).\tag{1}
\end{equation*}
其中 $\mathcal{N}(v)$ 是 $v$ 的邻居，$E(u,v)$ 是连接 $u,v$ 的边。\textbf{关键问题}就藏在这里：
传统 HGNN 里 $\mathrm{Propagate}(\cdot)$ 和 $\mathrm{Aggregate}(\cdot)$ 的参数\emph{依赖于具体的节点/边类型}，
所以模型被“锁死”在它训练过的那张图上，换一张图（类型集变了）就用不了——这正是 HiGPT 要打破的限制。

\section{五、数据集与任务}

\textbf{任务}：本文做的是\textbf{节点分类（node classification）}——给定异质图里的一个目标节点，
预测它属于哪一类。重点考察两种设定：
\begin{itemize}[leftmargin=2em]
  \item \textbf{少样本（few-shot）}：每类只给极少标注（1/3/5/10/20/40/60 个），在\emph{同一数据集}上测试。
  \item \textbf{零样本（zero-shot）/ 跨数据集迁移}：在数据集 A（IMDB）上训练，直接拿到数据集 B（DBLP/ACM）上测试，
        \textbf{完全不在 B 上训练}。这是检验“能否跨关系类型迁移”的关键。
\end{itemize}

\textbf{三个数据集}（统计见下表，源自原文 Table 4）：

\begin{table}[H]\centering
\caption{三个异质图数据集统计（数据源自原文 Table 4；带下划线者为分类的目标节点类型）。}
\begin{tabular}{llll}
\toprule
\textbf{数据集} & \textbf{节点（数量）} & \textbf{边类型（数量）} & \textbf{类别数} \\
\midrule
ACM  & \underline{论文}3025、作者5959、主题56、术语1902 & P-A 9949、P-S 3025、P-T 255619、P-P 5343 & 3 \\
DBLP & 作者4057、\underline{论文}14328、术语7723、会议20 & P-A 19645、P-C 14328、P-T 85810 & 4 \\
IMDB & \underline{电影}4278、导演2081、演员5257 & M-D 4278、M-A 12828 & 3 \\
\bottomrule
\end{tabular}
\end{table}

\begin{itemize}[leftmargin=2em]
  \item \textbf{IMDB}（电影库）：电影/导演/演员三种节点；电影分为 Action（动作）、Comedy（喜剧）、Drama（剧情）三类。
  \item \textbf{DBLP}（计算机文献库）：作者/论文/术语/会议四种节点；作者按 数据库/数据挖掘/人工智能/信息检索 四个研究领域分类。
  \item \textbf{ACM}（论文库）：论文/作者/主题三种节点；论文按 数据库/无线通信/数据挖掘 三类分类。
\end{itemize}

\textbf{评测指标}：Micro-F1（Mi-F1）、Macro-F1（Ma-F1）、AUC，三者\textbf{都是越高越好}。
为统一不同数据集的特征维度，所有节点都先用预训练的 Sentence-BERT 编码成同一维度的向量。

\begin{tip}{“IMDB 训练、DBLP 测试”为什么是硬核挑战}
IMDB 的节点是“电影/导演/演员”，DBLP 的节点是“作者/论文/会议”——\textbf{类型完全不重叠}，
关系也完全不同。传统模型在这种“跨域零样本”设定下基本等于瞎猜。HiGPT 能在这里显著领先，
正是它“在上下文里现学异质结构”能力的最好证明。
\end{tip}

\section{六、HiGPT 方法逐节精讲}

HiGPT 的整体架构见图~\ref{fig:arch}。一句话概括它的工作流：把异质图通过一个\textbf{“会现学”的
tokenizer} 编码成“图 token”，再把这些 token 连同人类指令一起喂给一个\textbf{冻结的 LLM（Vicuna）}，
让 LLM 直接输出预测。整条流水线有三大件：① 上下文异质图 tokenizer；② 两阶段异质感知指令微调；
③ 混合思维（MoT）指令增强。下面逐一拆解。

\begin{figure}[H]\centering
  \includegraphics[width=0.99\linewidth]{cand_p04_k00.pdf}
  \caption{HiGPT 总体架构（原文 Figure 1，图内英文为原文保留）。从左到右：
  输入异质图 + 异质图语料（IMDB/DBLP/ACM）$\to$
  \emph{语言增强的异质性表示}：用\textbf{冻结的} Sentence-BERT 把“这个节点是电影”“这条边是导演关系”
  等自然语言描述编码成类型向量 $T$ $\to$
  \emph{自适应参数化异质性投影器} + \textbf{HGT 层}（即上下文异质图 tokenizer，雪花=冻结、火焰=可训练）
  把图编码成\textbf{异质图 token}（按类型分组，type 1 / type 2 …）$\to$
  经\textbf{Graph Projector（图投影器）}对齐维度后，连同“带 \texttt{<Graph>} 占位符的人类指令”
  一起送入\textbf{冻结的 LLM（Vicuna）}。右侧三块分别对应：(a) 异质关系感知、(b) 同质关系感知、
  (c) 异质性感知微调。}
  \label{fig:arch}
\end{figure}

\subsection{第 1 件：上下文异质图 tokenizer（本文技术核心）}

这是攻克“关系类型异质性漂移”的关键。它的设计目标：\textbf{不把关系类型写死}，而是让模型
对\emph{任意}新图都能临时算出“该怎么编码”。它由两部分组成。

\textbf{(a) 用 HGT 做骨干。}\, tokenizer 可以用各种 HGNN 当骨干，本文以 HGT（异质图 Transformer）为例。
HGT 用“注意力 + 消息”来聚合不同类型邻居的信息（这是公式 (2)(3)，结构较繁，理解大意即可）：
\begin{equation*}
  \tilde h_v^{(l)}=\bigoplus_{\forall u\in\mathcal{N}(v)}
  \big(\mathrm{Attention}(u,e,v)\cdot\mathrm{Message}(u,e,v)\big),\qquad
  h_v^{(l)}=F_{\Theta_1}^{\tau(v)}\!\Big(\sigma\big(\tilde h_v^{(l)}\big)+h_v^{(l-1)}\Big).\tag{2}
\end{equation*}
这里 $\bigoplus$ 表示把所有邻居的贡献汇总，$\sigma(\cdot)$ 是激活函数，
$F_{\Theta_1}^{\tau(v)}$ 是一个\emph{依赖节点类型} $\tau(v)$ 的全连接层——\textbf{注意：正是这个“依赖类型”
让传统 HGT 换图就失效}。其中注意力与消息函数（$h$ 个注意力头）写作：
\begin{equation*}
  \mathrm{Attention}(u,e,v)=\operatorname*{Softmax}_{\forall u\in\mathcal{N}(v)}
  \Big(\!\!\underset{i\in[1,h]}{\big\|}\, F_{\Theta_2}^{\tau(u)}\!\big(h_u^{(l-1)}\big)\,
  W_1^{\rho(e)}\, F_{\Theta_3}^{\tau(v)}\!\big(h_v^{(l-1)}\big)^{\!\top}\Big),\quad
  \mathrm{Message}(u,e,v)=\underset{i\in[1,h]}{\big\|}\, F_{\Theta_4}^{\tau(u)}\!\big(h_u^{(l-1)}\big)\,W_2^{\rho(e)}.\tag{3}
\end{equation*}
$\|$ 表示把多个注意力头\textbf{拼接}起来，$W_1^{\rho(e)},W_2^{\rho(e)}$ 是\emph{依赖边类型} $\rho(e)$ 的矩阵。

\textbf{(b) 自适应参数化异质性投影器（关键创新）。}\, 既然“依赖类型的参数”是换图失效的根源，
作者干脆\textbf{不再为每种类型预先设定一套参数}，而是\textbf{临时“生成”这些参数}。具体地，
把上面那些依赖类型的参数 $F_{\Theta}^{\tau(v)}$、$W_i^{\rho(e)}$ 都由一个小网络 $F_{\Omega}$ 现场算出来：
\begin{equation*}
  \Theta_i=\{W_i^{\tau(v)};\,b_i^{\tau(v)}\}=F_{\Omega}\big(T_{\tau(v)}\big),\qquad
  W_i^{\rho(e)}=F_{\Omega}\big(T_{\rho(e)}\big).\tag{4}
\end{equation*}
其中 $F_{\Omega}$ 是一个参数为 $\Omega$ 的全连接层，$T_{\tau(v)}$、$T_{\rho(e)}$ 是“节点类型”和“边类型”
对应的\textbf{类型特征向量}。换句话说：\textbf{你给我任何一种新类型，我都能用同一个 $F_{\Omega}$ 当场算出
处理它所需的参数}——这就实现了“对任意异质图的现学现用”。

\begin{tip}{“现学”是怎么做到的？——把类型变成可计算的向量}
传统做法：图里有 5 种边，就预先准备 5 套参数；换成有 7 种边的新图，参数对不上，模型报废。\\
HiGPT 做法：不预存参数，而是把“边类型”本身编码成一个向量 $T_{\rho(e)}$，再用一个\emph{固定的}小网络
$F_{\Omega}$ 把这个向量\textbf{翻译}成该类型所需的参数。于是无论新图有几种、什么样的类型，
只要能给出它的类型向量，$F_{\Omega}$ 就能现场“变”出对应参数。这就是“in-context（上下文式）”的含义。
\end{tip}

\textbf{(c) 语言增强的异质性表示（Language-Enriched）。}\, 那“类型向量” $T_{\tau(v)}$ 从哪来？
答案是\textbf{用自然语言描述类型，再用语言模型编码}。比如把“电影”节点描述成
“This node represents a movie”，把（电影,\,to,\,导演）这条边描述成“The movie is directed by the director”。
为了让不同类型更可区分，每种类型用\emph{多句不同的话}来描述，再把编码结果取平均：
\begin{equation*}
  T_{\tau(v)}=\mathrm{Mean\text{-}Pooling}\big(\mathrm{Sentence\text{-}BERT}(\mathcal{S}_{\tau(v)})\big),\quad
  T_{\rho(e)}=\mathrm{Mean\text{-}Pooling}\big(\mathrm{Sentence\text{-}BERT}(\mathcal{S}_{\rho(e)})\big).\tag{5}
\end{equation*}
$\mathcal{S}_{\tau(v)}$、$\mathcal{S}_{\rho(e)}$ 是该类型的\textbf{一组语言描述}。例如那条“导演”边的描述集合是：
\begin{equation*}
  \mathcal{S}_{(\text{movie},\,\text{to},\,\text{director})}=\{\text{``The movie is directed by the director''},\
  \text{``The film features direction by the director''},\ \cdots\}.\tag{6}
\end{equation*}
\begin{tip}{为什么“用语言描述类型”能跨图迁移？}
因为语言是\textbf{所有图共享的“通用语”}。无论你的图来自电影库还是论文库，
“导演关系”和“作者关系”都能用英文句子描述，而 Sentence-BERT 对任何句子都给得出向量。
于是“类型”这个原本各图各异、无法对齐的东西，被统一映射进了同一个语义空间——
这就是 HiGPT 能把 IMDB 学到的本事迁移到 DBLP 的根本原因。
\end{tip}

\textbf{(d) 轻量文本—图对比对齐（预训练）。}\, 为了让 tokenizer 输出的“图向量”和 LLM 理解的“文本向量”
说同一种语言，作者借鉴 CLIP 的思路做了一次对比学习预训练：把图编码 $\hat H$ 和文本编码 $\hat T$ 归一化后，
让“配对的图—文本”相似、“不配对的”相离：
\begin{equation*}
  \hat H=\mathrm{norm}\big(\text{HG-Tokenizer}(X)\big),\quad
  \hat T=\mathrm{norm}\big(\text{LM-Tokenizer}(C)\big),\quad
  \mathcal{L}=\tfrac{1}{2}\big(\mathrm{CE}(\Lambda,y)+\mathrm{CE}(\Lambda^{\top},y)\big),\ \ \Lambda=(\hat H\hat T^{\top})\cdot\exp(\tau).\tag{7}
\end{equation*}
其中 $\mathrm{CE}(\cdot)$ 是交叉熵，$y=(0,1,\dots,n-1)^{\top}$ 是对比标签（即“第 $i$ 个图只和第 $i$ 段文本配对”），
$\tau$ 是温度系数。这一步把 tokenizer 平滑地“焊接”进了 HiGPT。

\subsection{第 2 件：两阶段异质感知指令微调（Heterogeneity-aware Fine-Tuning）}

有了图 token 还不够，得教 LLM \emph{看懂}它们。生成时，先把指令文本编码成 $X_I$，
再用一个投影器 $f_P$ 把图 token 对齐到文本同样的维度 $X_G=f_P(H)$，然后让 LLM 逐词生成答案 $X_O$：
\begin{equation*}
  p(X_O\mid X_G,X_I)=\prod_{i=1}^{L} p_{\Phi}\big(x_i\mid X_G,\,X_{I,<i},\,X_{O,<i}\big),\tag{8}
\end{equation*}
$\Phi$ 是 HiGPT 里可学习的参数。训练分\textbf{两个阶段}：

\textbf{阶段一：用异质图语料做指令预训练。}\, 目标是教会 LLM 区分“不同类型的图 token”。包含两个任务：
\begin{itemize}[leftmargin=2em]
  \item \textbf{异质关系感知（inter-type token matching，跨类型匹配）}：给 LLM 一串来自\emph{不同类型}的图 token，
        让它说出每个 token 是什么类型（如“1.电影、2.演员、3.导演”）。
  \item \textbf{同质关系感知（intra-type matching，同类型匹配）}：给 LLM 一串来自\emph{同一类型}的图 token
        和一组文字描述，让它把 token 和描述\textbf{正确配对/排序}。
\end{itemize}

\textbf{阶段二：异质性感知的下游微调。}\, 在阶段一之后，再用\emph{下游任务特定}的指令做有监督微调
（比如直接问“这部电影属于 action/comedy/drama 哪一类？”），把模型的推理能力定制到具体任务上。
表~\ref{fig:prompts} 给出三类任务的指令模板示例。

\begin{table}[H]\centering
\caption{异质图指令微调三类任务的提示模板示例（源自原文 Table 1，节选要点）。}
\begin{tabular}{@{}p{2.6cm}p{7.2cm}p{4.6cm}@{}}
\toprule
\textbf{任务} & \textbf{人类提问（Human Question）} & \textbf{HiGPT 回答} \\
\midrule
(a) 异质预训练 & 给定一个关于电影的异质图，有 3 类节点 \texttt{<DESC>}……不同类型的节点分别是：1.\texttt{<graph>}、2.\texttt{<graph>}……请依次给出各节点序列的类型。 & 根据图 token，类型应为：1.电影、2.……。 \\
\addlinespace
(b) 同质预训练 & 给定一个关于论文的异质图，有 4 类节点……“paper”节点为：\texttt{<graph>}；另给出论文的文字描述列表 \texttt{<DESC>}，请按图 token 顺序重排该列表。 & 图 token 与论文的匹配应为：\texttt{<ANSWER>}。 \\
\addlinespace
(c) 下游微调 & 给定一个关于电影的异质图……“movie”节点：\texttt{<graph>}（第 0 个为中心节点），“actor”节点：\texttt{<graph>}，“director”节点：\texttt{<graph>}。这部电影属于以下哪类：action/comedy/drama？ & 根据所给信息，该电影的可能类别是 Action。 \\
\bottomrule
\end{tabular}
\label{fig:prompts}
\end{table}

\begin{tip}{为什么要分两阶段？}
阶段一像“通识教育”：在海量异质图语料上，教模型\emph{认识图 token、分辨类型}——这是通用能力，
学一次终身受用。阶段二像“岗前培训”：只需\emph{极少}下游数据，就能把通用能力快速对齐到具体任务。
正因为阶段一打好了底子，阶段二哪怕只有“1 个样本”也能work——这正是 HiGPT 少样本超强的秘密。
\end{tip}

\subsection{第 3 件：混合思维（MoT）指令增强}

针对“数据稀缺”，作者让一个\textbf{强 LLM（GPT-3.5）自动\emph{扩写}指令}。具体用 7 种提示工程技巧
各生成一条带推理过程的指令，融合成“混合思维（Mixture-of-Thought）”：
\begin{itemize}[leftmargin=2em]
  \item \textbf{CoT（Chain-of-Thought，思维链）}：让模型给出中间推理步骤。
  \item \textbf{ToT（Tree-of-Thought，思维树）}：维护一棵“想法树”，系统地探索多条解题路径。
  \item \textbf{PanelGPT}：让多个模型像“专家组讨论”一样协作得出答案。
  \item \textbf{GKP（Generated Knowledge Prompting，生成式知识提示）}：先让模型生成相关背景知识，再据此作答。
\end{itemize}
一个巧妙之处：因为闭源大模型（如 ChatGPT）也可能答错、产生错误指导，作者把\textbf{正确答案（先验知识）
直接写进提示}，让 GPT-3.5 在“已知正确答案”的前提下\emph{倒推}出合理的中间推理步骤——
这样生成的推理过程既多样、又不会跑偏（见图~\ref{fig:mot}）。

\begin{figure}[H]\centering
  \includegraphics[width=0.92\linewidth]{cand_p05_k00.pdf}
  \caption{混合思维（MoT）指令增强（原文 Figure 2，图内英文为原文保留）：对同一个“人类提问 + 标准答案”，
  用 CoT/ToT/Panel/GKP 等多种方式各模拟出一段推理过程，从而把强 LLM 的多样推理能力“蒸馏”进 HiGPT，
  缓解数据稀缺。}
  \label{fig:mot}
\end{figure}

\section{七、实验与结果解读}

\textbf{实验设置}：基座 LLM 用开源的 Vicuna-7B；节点特征统一用 Sentence-BERT 编码；
对比三类基线——① 同质 GNN（SAGE、GAT）；② 异质 GNN（HAN、HGT、HetGNN）；
③ 异质图自监督方法（HGMAE、DMGI、HeCo）。HiGPT 有两个变体：\texttt{HiGPT-std}（直接给答案）
和 \texttt{HiGPT-cot}（带思维链）。围绕三个研究问题展开：RQ1 整体性能、RQ2 各组件消融、RQ3 图上下文学习。

\textbf{主结果（RQ1，源自原文 Table 2，节选关键行；Mi-F1，越高越好，加粗为最优）}：

\begin{table}[H]\centering
\caption{节点分类主结果对比（数值源自原文 Table 2 的 Micro-F1 行；“train-on/test-on”表示训练/测试数据集；
SSL 方法不做零样本故留空）。}
\begin{tabular}{llcccccc}
\toprule
\textbf{设定} & 训练$\to$测试 & GAT & HGT & HetGNN & HeCo & \textbf{HiGPT-std} & \textbf{HiGPT-cot} \\
\midrule
\multirow{2}{*}{少样本(同域)}
 & IMDB-1\ $\to$\ IMDB   & 0.4567 & 0.4977 & 0.4790 & 0.3874 & 0.5090 & \textbf{0.5360} \\
 & IMDB-5\ $\to$\ IMDB   & 0.5170 & 0.5003 & 0.5020 & 0.3385 & 0.6180 & \textbf{0.6320} \\
\midrule
\multirow{2}{*}{零样本(跨域)}
 & IMDB-1\ $\to$\ DBLP   & 0.1893 & 0.2573 & 0.2900 & -- & 0.3180 & \textbf{0.3500} \\
 & IMDB-1\ $\to$\ ACM    & 0.3567 & 0.3240 & 0.3743 & -- & 0.4160 & \textbf{0.4540} \\
\bottomrule
\end{tabular}
\end{table}

\textbf{怎么读这张表}：
\begin{itemize}[leftmargin=2em]
  \item \textbf{少样本碾压}：即使每类只给 1 个样本（IMDB-1），HiGPT-cot 就达到 0.536，
        把最强基线 HGT（0.498）甩开一大截；给到 5 个样本时更是 0.632 对 0.502。
  \item \textbf{零样本更见真章}：在“IMDB 训练、ACM 测试”这种\textbf{节点/关系类型完全不同}的跨域设定下，
        HiGPT-cot 达到 0.454，而所有传统 GNN 都在 0.32$\sim$0.37 附近徘徊（接近瞎猜）。
        这直接证明了 HiGPT 克服了“关系类型异质性漂移”。
  \item \textbf{CoT 普遍加分}：HiGPT-cot 几乎在每一行都优于 HiGPT-std，说明 MoT 注入的推理能力确实有效。
\end{itemize}

\textbf{消融实验（RQ2，源自原文 Table 3，IMDB 的 Mi-F1）}：逐个去掉关键组件看掉多少分。

\begin{table}[H]\centering
\caption{HiGPT 消融实验（数值源自原文 Table 3 的 Micro-F1；w/o S1=去掉阶段一指令预训练，
w/o HG=去掉异质图 token，w/o IA=去掉 MoT 指令增强）。}
\begin{tabular}{llcccc}
\toprule
\textbf{设定} & 训练$\to$测试 & w/o S1 & w/o HG & w/o IA & \textbf{HiGPT(完整)} \\
\midrule
\multirow{2}{*}{少样本}
 & IMDB-1\ $\to$\ IMDB & 失败(fail) & 0.3740 & 0.4260 & \textbf{0.5360} \\
 & IMDB-20\ $\to$\ IMDB & 失败(fail) & 0.5640 & 0.5620 & \textbf{0.6440} \\
\midrule
零样本
 & IMDB-1\ $\to$\ DBLP & 失败(fail) & 0.2980 & 0.2800 & \textbf{0.3500} \\
\bottomrule
\end{tabular}
\end{table}

\textbf{消融结论}：
\begin{itemize}[leftmargin=2em]
  \item \textbf{去掉阶段一预训练（w/o S1）直接“失败”}：只用下游少量数据微调，模型在 15/50/100 个 epoch 下
        都\emph{无法给出完整准确的回答}。这说明“在大规模异质图语料上的通识预训练”是 HiGPT 的\textbf{命脉}。
  \item \textbf{去掉图 token（w/o HG）掉分严重}：1-shot 时从 0.536 跌到 0.374。说明图结构信息（而不只是文本）
        是性能的关键来源，尤其在样本稀疏时。
  \item \textbf{去掉 MoT 增强（w/o IA）也明显掉分}：印证了 MoT 在缓解数据稀缺上的价值。
\end{itemize}

\textbf{图上下文学习（RQ3，源自原文 Figure 3）}：最惊艳的发现是——\textbf{“1 个样本 + 图上下文示例”能打败
“60 个样本但无示例”}。即在推理时\emph{不更新任何参数}，只在提示里塞进几个“图问答示例”，
模型性能就能超过用 60 个样本训练的版本。更有意思的是：在 ACM 上测试时，
\textbf{放入“不相关的 DBLP 示例”反而效果最好}——说明 HiGPT 真正学到的是\emph{从图结构中提取信息}的通用能力，
而非死记某个数据集，因而展现出极强的迁移性。

\section{八、小白知识点卡片（速查）}
\begin{itemize}[leftmargin=2em]
  \item \textbf{异质图（heterogeneous graph）}：节点和边\emph{有多种类型}的图（学术网络、推荐系统等）；反之为同质图。
  \item \textbf{HGNN}：异质图神经网络，按节点/边类型分别设计消息传递与聚合。
  \item \textbf{元关系 / 元路径（meta relation / metapath）}：描述异质图里“类型层面”关系的三元组及其串联。
  \item \textbf{关系类型异质性漂移}：不同图的关系\emph{种类、含义}都不同，导致旧模型无法迁移到新图。
  \item \textbf{图指令微调}：把图夹进“指令 + 回答”里训练 LLM，让它边看图边听话边作答。
  \item \textbf{上下文 tokenizer（in-context）}：不预存类型参数，靠 $F_{\Omega}$ 现场把“类型向量”翻译成参数，从而适配任意新图。
  \item \textbf{语言增强表示}：用自然语言描述节点/边类型并用 Sentence-BERT 编码，借“通用语言”实现跨图对齐。
  \item \textbf{MoT（混合思维）}：用 CoT/ToT/PanelGPT/GKP 多种提示让强 LLM 自动扩充带推理的指令。
  \item \textbf{CoT / ToT / GKP}：思维链 / 思维树 / 生成式知识提示，三种引导 LLM 推理的提示技巧。
  \item \textbf{零样本 / 少样本（zero-/few-shot）}：完全不在目标数据上训练 / 仅用极少样本训练。
  \item \textbf{Micro-F1 / Macro-F1 / AUC}：分类任务常用的三种精度指标，均越高越好。
  \item \textbf{图上下文学习（Graph ICL）}：推理时只在提示里给几个“图问答示例”，不更新参数就能提升表现。
\end{itemize}

\section{九、一句话总结与谱系定位}

\textbf{一句话}：HiGPT 把“图指令微调”从同质图推广到\textbf{异质图}——核心是一个\emph{带语言增强、能在上下文里
现学、适配任意关系类型}的 in-context 异质图 tokenizer（靠 $F_{\Omega}$ 把“类型”当场翻译成参数），
再配上\textbf{两阶段异质感知指令微调}和\textbf{MoT 指令自动扩充}，
从而做出一个\emph{不用在下游数据上微调}、就能在少样本/零样本与跨数据集迁移上都很强的“通用异质图大模型”。

\textbf{在“图 $\times$ 大模型”谱系里的位置}：HiGPT 出自香港大学 HKUDS 团队，
是其 \textbf{GraphGPT}（同质图指令微调）的\textbf{异质图升级版}——
GraphGPT 解决“让 LLM 读懂\emph{同质}图”，HiGPT 进一步攻克“\emph{异质}图的类型漂移与跨域迁移”。
与本系列里其他路线对照：G-Retriever 走“图上检索增强（GraphRAG）”、LLaGA 走“投影器对齐”、
InstructGraph/GraphWiz 走“指令微调 + DPO”——你可以对照着看它们\emph{各自把“图的信息”塞进 LLM 的方式有何不同}，
而 HiGPT 的独到之处是把火力集中在了“异质性的可迁移编码”上。

\textbf{动手复现}：官方代码 \url{https://github.com/HKUDS/HiGPT}（含异质图 tokenizer、
两阶段指令数据构造、MoT 增强脚本与 Vicuna 训练流程），项目主页 \url{https://HiGPT-HKU.github.io}，
是入门“异质图 $\times$ 大模型”很好的练手项目。

\end{document}
